\begin{center}
  \section{\capitalisewords{Artificial Intelligence in Everyday Life}}
  Angana Mukhopadhyay, 3$^{rd}$ Year, ECE
\end{center}

\setcounter{figure}{0}
\setcounter{table}{0}

\begin{multicols}{2}
  \begin{abstract}
    \bfseries
    Artificial Intelligence (AI) has rapidly evolved from a futuristic idea into an essential part of everyday life, influencing how people communicate, work, travel, shop, and learn. From voice assistants such as Siri, Alexa, Cortana, Bixby, and Google Assistant to personalized recommendations on Netflix and online shopping platforms, AI is quietly transforming daily experiences through smart automation and data-driven decision-making. In healthcare, AI helps doctors detect diseases faster and improve patient care, while in banking and cybersecurity it enhances fraud detection and digital security. Smart homes, self-driving technologies, and AI-powered educational tools further demonstrate how intelligent systems improve convenience, efficiency, and productivity.

    \noindent
    This report explores the growing role of Artificial Intelligence in modern society and examines its practical applications across multiple sectors. It highlights benefits such as time-saving automation, improved accuracy, and enhanced user experiences, while also discussing privacy concerns, ethical issues, job displacement, and overdependence on technology. By examining real-world applications and technologies such as edge computing and Natural Language Processing, the report shows how AI is reshaping human life. As Artificial Intelligence continues to advance, understanding its impact becomes increasingly important for building a balanced and responsible technological society.

    \vspace{0.18em}
    {\justifying\noindent\textbf{Keywords:} Artificial Intelligence, Machine Learning, Natural Language Processing, Computer Vision, Robotics, Smart Homes, Healthcare, Cybersecurity, Automation.\par}
  \end{abstract}
  \vspace{-0.6em}
  \subsection*{\centering\uppercase{Introduction}}
  Artificial Intelligence has emerged as one of the most revolutionary technologies of the twenty-first century, transforming the world at an extraordinary pace. What once existed mainly in science-fiction films and futuristic predictions is now deeply embedded in everyday human life. From voice assistants answering questions within seconds to smart applications predicting user preferences, AI has become an invisible intelligence powering modern society. In healthcare, education, banking, transportation, entertainment, and communication, AI is continuously changing how people think, interact, and make decisions.

  \noindent
  The increasing dependence on intelligent systems has created a digital era in which machines can learn from data, recognize patterns, solve complex problems, and mimic aspects of human behaviour. Technologies such as facial recognition, self-driving vehicles, chatbots, recommendation algorithms, and smart-home devices demonstrate how AI is making life more convenient, efficient, and connected.

  \noindent
  At the same time, the rapid advancement of Artificial Intelligence has raised important concerns related to privacy, cybersecurity, ethical decision-making, employment opportunities, and human dependence on machines. AI is no longer only a technological innovation; it has become a powerful force shaping the future of humanity. Understanding its impact on everyday life is essential to ensure that technological progress benefits society while maintaining ethical responsibility and human values.

  \subsection*{\centering\uppercase{Meaning of Artificial Intelligence}}
  Artificial Intelligence refers to the simulation of human intelligence in machines that are programmed to think, learn, and make decisions. It includes technologies such as Machine Learning, Natural Language Processing, Computer Vision, and robotics. AI systems analyze large amounts of data to identify patterns, make predictions, and improve their performance over time without being explicitly reprogrammed for every task.

  \noindent
  AI is widely used in healthcare, finance, transportation, education, and entertainment to automate tasks and enhance efficiency. It can range from simple rule-based systems to advanced neural networks that imitate selected functions of the human brain, enabling machines to perform complex problem-solving and adapt to new situations.

  \subsection*{\centering\uppercase{How Everyday AI Systems Function}}
  The basic operation of AI systems can be described through the following stages:

  \begin{itemize}\justifying
    \item \textbf{Data Collection:} AI systems gather data from users, devices, sensors, websites, and applications. The data may include text, images, voice commands, location information, or user behaviour.
    \item \textbf{Data Processing:} The collected information is cleaned and organized so that the AI system can analyze it effectively. Algorithms process the data to identify useful patterns and relationships.
    \item \textbf{Decision-Making:} Based on learned patterns, the system makes predictions, recommendations, or decisions. Voice assistants answer questions, streaming platforms recommend films, and navigation applications suggest faster routes.
    \item \textbf{Machine Learning:} AI systems use Machine Learning techniques to learn from historical data and improve their performance over time without being explicitly programmed for every situation.
  \end{itemize}

  \subsection*{\centering\uppercase{Major Types of Artificial Intelligence}}
  \subsubsection*{Narrow AI}
  Narrow AI is designed to perform a specific task. It works within a limited range and cannot think like a human. Its main features are task-specific operation, speed, and accuracy. It is the most common form of AI currently used in everyday systems. Examples include Google Assistant, Netflix recommendations, and Google Maps.

  \subsubsection*{General AI}
  General AI, also known as Strong AI, describes a system that can think, learn, and perform multiple tasks like a human being. It would be able to understand, reason, solve problems independently, learn automatically, and adapt its knowledge across different tasks.

  \subsubsection*{Super AI}
  Super AI refers to an advanced form of Artificial Intelligence that could surpass human intelligence in thinking, decision-making, creativity, and problem-solving. Such a system would be capable of independent learning and handling highly complex problems.

  \subsection*{\centering\uppercase{AI-Powered Technologies Around Us}}
  \subsubsection*{AI in Healthcare and Medicine}
  Artificial Intelligence is used in healthcare to improve diagnosis, treatment, and patient care. It helps with early disease detection, assists doctors in medical diagnosis, analyzes medical images such as X-rays and MRI scans, supports drug discovery and research, and enables remote patient monitoring. Robotic-assisted surgery and hospital systems that analyze patient information are important examples.

  \subsubsection*{AI in Education and E-Learning}
  AI is transforming education by making learning more personalized, interactive, and accessible. It supports personalized learning based on student performance, smart tutoring systems, virtual teachers, automated grading, language translation, and speech-to-text tools. Adaptive learning applications can adjust their difficulty level according to the progress of the learner.

  \subsubsection*{AI in Navigation and Transportation}
  Artificial Intelligence is used in navigation and transportation for route optimization, traffic prediction, ride-sharing services, maps, and autonomous vehicles. It improves travel efficiency, reduces traffic congestion, enhances road safety, and provides real-time updates, making transportation faster and more reliable.

  \subsubsection*{AI in E-Commerce}
  E-commerce platforms use AI to improve the shopping experience through product recommendations, personalized advertisements, chatbots, and demand forecasting. AI analyzes user behaviour to suggest relevant products, optimize pricing, and detect fraud. This makes online shopping faster, more convenient, and more personalized for customers and businesses.

  \subsubsection*{AI in Social Media, Banking, Smartphones, and Smart Homes}
  In social-media platforms, AI analyzes user behaviour to recommend posts, reels, videos, friends, and targeted advertisements. In banking and digital payments, it detects fraud, monitors transactions, supports credit scoring, and provides continuous customer support through chatbots.

  \noindent
  In smartphones and applications, AI powers voice assistants, facial recognition, predictive text, camera enhancement, and personalized app suggestions. In smart homes and security systems, AI automates lights, thermostats, and appliances while monitoring security cameras and detecting unusual activity. These applications improve convenience, efficiency, personalization, and security.

  \subsubsection*{AI in Robotics}
  {\justifying Artificial Intelligence is making robots smarter, faster, and more efficient. Intelligent robots can learn, recognize objects, interact with their surroundings, and make decisions. Self-driving cars, robotic surgery, smart home assistants, and industrial automation demonstrate how robots improve accuracy, reduce human effort, and perform dangerous tasks safely. Machine Learning, Computer Vision, and sensors help them respond intelligently to their environment.\par}

  \subsubsection*{Cancer Detection and Treatment Support}
  Cancer treatment involves methods such as surgery, chemotherapy, radiation therapy, immunotherapy, and targeted therapy. Surgery removes tumours, while chemotherapy and radiation destroy harmful cells. Immunotherapy strengthens the body's immune system, and targeted therapy acts on specific cancer mechanisms.

  \noindent
  AI supports cancer care by helping to detect disease early through medical-image analysis and data prediction. Early diagnosis increases the possibility of successful treatment and recovery. Regular medical check-ups, a healthy lifestyle, and modern healthcare technologies play important roles in prevention and treatment. AI-powered MRI scanners and AI-based mammography systems are two technologies used to assist cancer detection.

  \subsubsection*{AI in Cybersecurity}
  AI-based cybersecurity tools can detect hackers, malware, phishing attacks, and suspicious activity in real time before major damage occurs. By continuously analyzing large quantities of data, these systems learn attack patterns and respond automatically to cyber threats. Biometric authentication, AI firewalls, and intelligent antivirus software help protect personal data, banking systems, businesses, and online platforms.

  \subsection*{\centering\uppercase{Challenges and Risks of Artificial Intelligence}}
  \begin{itemize}\justifying
    \item \textbf{Privacy and Data Security:} AI systems collect and analyze personal information such as photographs, locations, banking details, and online activities. Weak protection can lead to data theft, privacy violations, and cybercrime.
    \item \textbf{Job Displacement:} AI-powered automation can replace human workers in repetitive jobs such as manufacturing, customer support, and data entry. Workers may need to learn new skills as employment requirements change.
    \item \textbf{Cybersecurity Threats:} AI can be misused to create advanced cyberattacks, deepfake videos, phishing scams, and automated hacking tools that are difficult to detect.
    \item \textbf{Bias and Discrimination:} If training data contains bias, an AI system may produce unfair decisions in hiring, loan approval, or law enforcement.
    \item \textbf{Lack of Transparency:} Some systems operate like black boxes, making it difficult even for developers to explain how decisions are produced. This reduces trust in critical applications such as healthcare and finance.
    \item \textbf{Threat to Human Control:} Highly advanced systems may become difficult to control if they make independent decisions without proper monitoring.
    \item \textbf{System Errors and Failures:} Incorrect predictions and software failures can cause serious problems in healthcare, transportation, and defence systems.
    \item \textbf{High Development Cost:} Advanced AI requires powerful computers, large datasets, skilled experts, and continuous maintenance, making it expensive for smaller organizations.
  \end{itemize}

  \subsection*{\centering\uppercase{The Future of Artificial Intelligence}}
  The future applications of Artificial Intelligence include:

  \begin{itemize}\justifying
    \item \textbf{Smarter Healthcare:} earlier disease detection, improved diagnosis, robotic surgery, and personalized treatment plans.
    \item \textbf{Self-Driving Vehicles:} autonomous cars, buses, and trucks designed to improve safety and travel efficiency.
    \item \textbf{Advanced Robotics:} intelligent machines performing complex tasks in industries, hospitals, homes, and space exploration.
    \item \textbf{Personalized Education:} learning systems that adapt to the speed and requirements of individual students.
    \item \textbf{Improved Cybersecurity:} faster detection of cyber threats, hacking attempts, and online fraud.
    \item \textbf{Smart Homes and Cities:} automated lighting, security, energy management, traffic control, and public services.
    \item \textbf{Scientific Research:} analysis of complex data to accelerate discoveries in medicine, climate studies, and space research.
    \item \textbf{Human--AI Collaboration:} systems working alongside people to improve efficiency, creativity, and problem-solving.
    \item \textbf{Business Automation:} automation of repetitive tasks, customer analysis, decision support, and productivity improvement.
  \end{itemize}

  \subsection*{\centering\uppercase{Conclusion}}
  Artificial Intelligence has established itself as a transformative technology in many areas of everyday life. It has improved efficiency, personalized experiences, and created new opportunities for innovation in healthcare, banking, education, social media, transportation, and other sectors. However, its rapid integration also creates legal, ethical, and technological challenges.

  \noindent
  AI has a dual nature: it offers enormous potential while also introducing serious risks. Concerns related to bias, data privacy, job displacement, security, and excessive dependence on intelligent systems must be addressed as the technology develops. Responsible development and use are essential to ensure that the benefits of AI are distributed fairly and its risks are minimized.

  \noindent
  The rise of Artificial Intelligence is therefore not only a story of innovation and convenience; it is also a challenge that demands responsibility, ethics, and human awareness. As machines become more capable of learning and decision-making, society must ensure that human values, creativity, and control remain at the centre of technological progress.

  \subsection*{\centering\uppercase{References}}
  \begin{justify}
  \footnotesize
  \begin{enumerate}[leftmargin=*,itemsep=0.20em]
    \item J. Kaplan and M. Haenlein, “Siri, Siri, in My Hand: Who's the Fairest in the Land? On the Interpretations, Illustrations, and Implications of Artificial Intelligence,” \textit{Business Horizons}, vol. 62, no. 1, pp. 15--25, 2019.
    \item S. Russell and P. Norvig, \textit{Artificial Intelligence: A Modern Approach}, 4th ed. Hoboken, NJ, USA: Pearson, 2021.
    \item D. S. Schiff, A. Ayesh, L. Musikanski, and J. C. Havens, “IEEE 7010: A New Standard for Assessing the Well-Being Implications of Artificial Intelligence,” \textit{IEEE Technology and Society Magazine}, vol. 39, no. 3, pp. 50--55, 2020.
    \item M. Kasinidou, S. Kleanthous, M. Busso, M. Rodas, J. Otterbacher, and F. Giunchiglia, “Artificial Intelligence in Everyday Life 2.0: Educating University Students from Different Majors,” 2024.
    \item R. Luckin, \textit{Machine Learning and Human Intelligence: The Future of Education for the 21st Century}. London, U.K.: UCL Institute of Education Press, 2018.
    \item B. Marr, \textit{Artificial Intelligence in Practice: How 50 Successful Companies Used AI and Machine Learning to Solve Problems}. Hoboken, NJ, USA: Wiley, 2019.
    \item G. Tecuci, “Artificial Intelligence,” \textit{Wiley Interdisciplinary Reviews: Computational Statistics}, vol. 4, no. 2, pp. 168--180, 2012.
    \item IEEE, “Artificial Intelligence and Machine Learning Applications,” 2024.
  \end{enumerate}
  \end{justify}
\end{multicols}
